\section{Discussion}\label{sec:discussion}

\paragraph{Why the multiplicative form.} The additive alternative keeps the same covariates and tab effects, but replaces the product of user examination and item attractiveness by an additive logit. Against the full DAX arm it is lower by $+6.09$ [0.96] (per-user t test, Holm-adjusted $p=5.43\times10^{-35}$), and it keeps 59.7\% of the full gain (Table~\ref{tab:ablate}). Under the generative model of the method section, a non-click is weak evidence against an item for a user who rarely clicks anything, and weak evidence in a scenario that rarely yields clicks. The multiplicative form encodes that dependence; the additive form, once the user effect is absorbed, treats the two cases alike. The comparison is a fit comparison under one protocol, and it does not show that the multiplicative form is correct in an absolute sense.

\paragraph{What the item-level part contributes.} Permuting $\alpha$ within tags keeps 41.6\% of the gain, while permuting it across items removes all of it (Table~\ref{tab:ablate}). A model that used only a tag mean would therefore capture part of the signal, and the remainder is carried by item deviations from that mean. This is the reason the item-level estimate is retained, but the decomposition is only as strong as the permutation controls, which keep the marginal distribution of $\alpha$ fixed.

\paragraph{Why the lifelong channels are not interchangeable with the item channels.} Shuffling the lifelong tables across users removes the LEAF gain, and removing all lifelong channels removes $+17.88$ [2.04] (Table~\ref{tab:leafctl}). The policy channel contributes less ($+5.01$ [1.20]) and its split-level test is not significant after Holm adjustment. The signal that transfers to randomly exposed items is therefore mainly the users' own lifelong affinities for content keys, and the in-pool policy channel adds less on top of it.

\paragraph{Complementarity rather than stand-alone accuracy.} DAX is used as an increment on top of a backbone, and the tables report only that increment. The increment is positive on every backbone with a result (Table~\ref{tab:backbone}), and the combination is above the backbones on which it is added. The paper does not report DAX as a stand-alone ranker: the claim is that its score carries information that the backbone scores do not, under the post-hoc rule.

\paragraph{Lifelong deep models.} DIN with DAX is the closest comparator in the main family. LEAF-LGN with DAX is above it by $+5.73$ [1.64] with a Holm-adjusted Nadeau-Bengio p of $4.56\times10^{-2}$, which ties with the QP family as the weakest split-level margin of the seven comparators. The lifelong deep baselines are tuned over four configurations each. Whether SIM with DAX, once run, changes the ordering is open, and the paper makes no claim about it.

\paragraph{What the remaining headroom would require.} The gains are concentrated in the item-level appeal estimate and in the lifelong user channels. Further gains would need a different source of information about item quality for random viewers, such as richer content for items or a direct measurement of uniform-exposure responses, and not a larger search over the present hyperparameters. This is an interpretation of the controls and not a tested claim.
